% page 542 of the facsimile (image p-541 of the scan)
% block 160.0 x 233.3 mm ; left margin 8.0 mm
\pgc{-12.700mm}{0.000mm}{
\l{\cel{3}534}
\l{}
\l{\sou{skafandro}. (\sou{tekn.}) Aparato aptigita pri kontenar plunjisto}
\l{\cel{3}qua, per lu, povas laborar sub la aquo-surfaco, per res-}
\l{\cel{3}pirar la aero quan sendas a lu pumpilo situita exter la}
\l{\cel{3}aquo. - FIS.}
\l{}
\l{\sou{skalo}. I. Eskalero portebla, ek du montanti, an qui adjuste-}
\l{\cel{3}sis traversi. - II. (\sou{metaf.})\cel{2}(a) Serio acensanta o decen-}
\l{\cel{3}santa (kom ex.: la skalo de la enti, la skalo sociala, la}
\l{\cel{3}skalo diatonika, la skalo de la kolori); (b) Ligno, dividi-}
\l{\cel{3}ta aden parti (gradi), e c., uzebla por mezurar.- DEFIRS.}
\l{}
\l{\sou{skalaro}. (\sou{konkologio}). Genero de moluski karnivora, gasteropo-}
\l{\cel{3}da, prosobrankia, tipo di la familio \textquotedbl{}skalaridei\textquotedbl{} - kunros-}
\l{\cel{3}tro retrotirebla, palpili longa, konko turmatra, solida,}
\l{\cel{3}briloza, kun multa spiri kostoza. - DEFIRS.}
\l{}
\l{skaleno. (geom.) Triangulo di qua la tri lateri esas inter-}
\l{\cel{3}ne-egala. -}
\l{}
\l{skalio. (\sou{zool.}) I. Krusto qua kovras la korpo di la ostri o di}
\l{\cel{3}la tortugi. II. Kuraso qua kovras la korpo di la tortugi.}
\l{}
\l{}
\l{\sou{skalpar}. (\sou{trans.})\cel{2}Deprenar la pelo di la kranio. - DEFIR.}
\l{}
\l{\sou{skalpelo}. (\sou{kirurg.)}\cel{2}Instrumento kun lamo fixa, pintoza, kun}
\l{\cel{3}un o du aristi, quan onu uzas por dissekar. - DEFIRS.}
\l{}
\l{\sou{skamoneo}. Gumo rezina, purgiva, quan onu extraktas de la ra-}
\l{\cel{3}diko di speci diversa de konvolvulo mikra. - DEFIS.}
\l{}
\l{\sou{skandar}. (\sou{trans.}) I. Pronuncar (verso) acentizante forte, o}
\l{\cel{3}deskompozante, ica pos ita, singla ek la pedi (metrala o}
\l{\cel{3}silabala) ye qui lu konsistas. - II. Kantar, plear (frazo}
\l{\cel{3}muzikala) acentizante forte la tempi di singla mezuro. -}
\l{\cel{3}DEFIRS.}
\l{}
\l{\sou{skandalar}. (\sou{trans., per)}. Igar sua kunhomi en danjero di etiko-}
\l{\cel{3}falio, di moro-falio, per lia tenteso imitar onua mala exem-}
\l{\cel{3}plo. - II. Shokar psikale per la prizento regretinda di mala}
\l{\cel{3}exemplo. - DEFIRS.}
\l{}
\l{\sou{skapulo}. (\sou{anat.}) Osto triangula, larja e plata, ye qua konsis-}
\l{\cel{3}tas parto dopa di la shultro, e kun qua la osto di la bra-}
\l{\cel{3}kio formacas artiko.- EL.}
\l{}
\l{\sou{skapulario.(religio kristan.}) I. Parto di la vesto di ula re-}
\l{\cel{3}gulbri, qua kovras la robo avane e dope. - II. Asembluro}
\l{\cel{3}ek du mikra peci de drapo sakrigita, sur qui esas brodita}
\l{\cel{3}portreto de la santa Virgino (kelke simila ye la skapula-}
\l{\cel{3}rio di la regulieri) quan ula personi +weras sub sua vesti.}
\l{\cel{3}- DEFIRS.}
}
